\subsection{Crucial Features}
\noindent\textbf{Chunking along space.}
The computational domain is partitioned exclusively along its spatial dimensions, while the complete
temporal sequence associated with each spatial point is preserved.
This constraint ensures that event continuity, persistence, and sub-daily temporal structure remain
available throughout the calculation.
\noindent\textbf{Batching and out-of-memory recovery.}
Spatial batches are dynamically adapted according to the GPU memory available at execution time,
allowing the same workflow to operate across devices with different VRAM capacities.
If an out-of-memory condition occurs, the batch size is reduced and the affected spatial subset is
retried without truncating or partitioning its temporal record.

\noindent\textbf{Conversion to single precision.}
Single-precision floating-point representation reduces both memory requirements and computational cost
relative to double precision, allowing larger spatial batches to be processed on the GPU.
The numerical impact of this conversion is explicitly evaluated against a double-precision CPU
reference to quantify any differences introduced by reduced precision.



\noindent\textbf{Preserved consistency.}
GPU-derived results are systematically compared with the reference CPU implementation to verify that
computational acceleration does not materially alter the scientific outcome.
Consistency is evaluated for intermediate fields, threshold-exceedance classifications, and derived
quantities, with particular attention to differences close to the detection threshold.

\noindent\textbf{Improved performance.}
GPU execution substantially reduces the computational time required to process large hourly SST
datasets compared with the corresponding CPU implementation.
Performance is quantified through execution time, speedup, throughput, memory consumption, and the
relative contributions of host-to-device transfer, GPU computation, and device-to-host transfer.


